% Documento mínimo de teste do ambiente.
% Compilação: latexmk -lualatex main.tex
\documentclass[a4paper,12pt]{article}

\usepackage{fontspec}
\usepackage[brazilian]{babel}
\usepackage{amsmath,amssymb,mathtools,amsthm}
\usepackage{graphicx}
\usepackage{pgfplots}
\pgfplotsset{compat=1.18}
\usepackage{microtype}
\usepackage{csquotes}
\usepackage[backend=biber,style=alphabetic]{biblatex}
\addbibresource{refs.bib}
\usepackage[hidelinks]{hyperref}

\begin{document}

\begin{center}
  {\LARGE Teste de configuração do ambiente\par}
\end{center}

\section{Introdução}

Este parágrafo verifica a acentuação do português: ação, coração, pé, avô,
à noite, ênfase, órbita, pão, põe, fácil, cálculo, função e série.
Também verifica aspas: \enquote{citação curta}.

A série de Basel converge para
\begin{equation}\label{eq:basel}
  \sum_{n=1}^{\infty} \frac{1}{n^{2}} = \frac{\pi^{2}}{6},
\end{equation}
resultado que admite uma demonstração elementar por integral dupla
\autocite{Apostol1983proof}. A \autoref{fig:somas} mostra as somas parciais
$S_N = \sum_{n=1}^{N} 1/n^{2}$ aproximando-se do valor
em~\eqref{eq:basel}.

\begin{figure}[htbp]
  \centering
  % Somas parciais S_N = \sum_{n=1}^N 1/n^2 (valores calculados e conferidos;
% 5 casas decimais) e o limite \pi^2/6 \approx 1,64493.
\begin{tikzpicture}
  \begin{axis}[
      width=0.82\linewidth, height=0.45\linewidth,
      xmin=0, xmax=21, ymin=0.9, ymax=1.76,
      xlabel={$N$}, ylabel={$S_N$},
      xtick={1,5,10,15,20},
      ytick={1,1.2,1.4,1.6},
      yticklabel style={/pgf/number format/.cd, use comma, fixed, precision=1},
      axis lines=left, tick align=outside,
      every axis plot/.append style={thick},
    ]
    \addplot[only marks, mark=*, mark size=1.6pt] coordinates {
      (1,1.00000) (2,1.25000) (3,1.36111) (4,1.42361) (5,1.46361)
      (6,1.49139) (7,1.51180) (8,1.52742) (9,1.53977) (10,1.54977)
      (11,1.55803) (12,1.56498) (13,1.57089) (14,1.57600) (15,1.58044)
      (16,1.58435) (17,1.58781) (18,1.59089) (19,1.59366) (20,1.59616)
    };
    \addplot[dashed, domain=0:21] {pi^2/6}
      node[pos=0.04, above right] {$\pi^2/6$};
  \end{axis}
\end{tikzpicture}

  \caption{Somas parciais $S_N$ da série~\eqref{eq:basel} para
    $1 \le N \le 20$; a reta tracejada indica o limite $\pi^{2}/6$.}
  \label{fig:somas}
\end{figure}

\printbibliography

\end{document}
